\documentclass[aps,prd,a4paper,10pt, twocolumn,amsmath,showpacs,superscriptaddress,nofootinbib,preprintnumbers,notitlepage]{revtex4-2}
\usepackage{dcolumn}
\usepackage{epstopdf}
\usepackage{bm}
\usepackage{braket}
\usepackage{enumitem}
\usepackage{soul}
\usepackage{pifont}
\usepackage[labelfont= bf, skip=5pt, font=small]{caption}
\usepackage{graphicx, wrapfig, subcaption}
\usepackage{array, tabularx, tabulary, booktabs, multirow}
\usepackage{textcomp, cmap, comment, epigraph}
\usepackage[dvipsnames]{xcolor}
\usepackage{orcidlink}
\usepackage[utf8]{inputenc}
\usepackage[spanish, es-tabla]{babel}
\usepackage[version=4]{mhchem}
\usepackage{hyperref}
\usepackage{geometry}
\usepackage{amsfonts, amssymb, amsthm, mathtools} 
\usepackage{amsmath}
\usepackage{float}
\usepackage{placeins}
\usepackage{siunitx}

\geometry{left=15mm, right=15mm, top=16mm, bottom=18mm}
\DeclareMathSizes{12}{16}{12}{10}
\selectlanguage{spanish}

\makeatletter
\newcommand{\fmarki}{*}
\newcommand{\fmarkii}{\ensuremath{\dagger}}
\newcommand{\fmarkiii}{\ensuremath{\ddagger}}  
\def\@fnsymbol#1{{\ifcase#1\or \fmarki\or \fmarkii\or \fmarkiii\else\@ctrerr\fi}}
\makeatother

\renewcommand{\thesection}{\Roman{section}}
\renewcommand{\abstractname}{}

\addto\captionsspanish{\renewcommand{\abstractname}{}}
\def\restrict#1{\raise-.5ex\hbox{\ensuremath|}_{#1}}
\def\sen{\mathop{\mbox{\normalfont sen}}\nolimits}
\graphicspath{ {images/} }
\hypersetup{colorlinks=true,linkcolor=MidnightBlue,citecolor=OliveGreen,urlcolor=RoyalBlue,pdftitle={},pdfauthor={Nicolás Clotta}}

\begin{document}

\title{Cuaderno de Laboratorio 6}

\author{Nicolás Clotta}
\email{nicolas.clotta@gmail.com}
\affiliation{Departamento de Física, Facultad de Ciencias Exactas y Naturales, Universidad de Buenos Aires.}

\date[Fechado: ]{\today}
%\begin{abstract}
%    \noindent
%\end{abstract}

\maketitle

\section{Minor loops, paredes de dominio}

La posibilidad de medir la magnetización 

\section{Lake Shore VSM}
\begin{itemize}
\item $m_s$ momento de saturación: A partir de valores del campo entre los cuales el moment magnético es lineal, 

\item $H_{ci}$ coercitividad intrínseca

\item $H_c$ coercitividad

\item $m_r$ remanencia

\item $B_r$ inducción residual

\item $\displaystyle\frac{m_r}{m_s}$ squareness ratio

\item $B\times H_{max}$ max. energy product
\end{itemize}

\bibliographystyle{plain}
\bibliography{cuaderno6}




\end{document}
